\centerline{\bf \Huge II Тренировочная Олимпиада}

\begin{flushright}

\end{flushright}


\problem{\textbf{1}} Помогите Мише расположить по кругу числа 1, 2, 3, 4, 5, 6, 7, 8 так,
чтобы каждое число было кратно разности двух соседних с ним чисел. 

\problem{\textbf{2}} В строку записаны 5 чисел. Могло ли оказаться так, что сумма любых двух
соседних чисел – отрицательна, а сумма всех пяти чисел – положительна?

\problem{\textbf{3}} На доске выписаны числа от 1 до 2150. Каждую минуту каждое число подвергается следующей
операции: если число делится на 100, то его делят на 100; если же не делится, то из него
вычитают 1. Найдите наибольшее среди чисел на доске через 81 минуту.


\problem{\textbf{4}} В шахматном турнире на кубок клуба четырех коней участвовали 2018 шахматистов, причем
каждый сыграл с каждым ровно один раз. Остап Бендер рассказал, что по итогам турнира
каждый выиграл ровно столько матчей, сколько сыграл вничью. Не врет ли Остап? (победа - 2 очка, ничья - 1, поражение - 0)


\problem{\textbf{5}} Знайка утверждает, что существует 8-значное число, которое при делении на свою первую
цифру дает в остатке 1, при делении на вторую цифру дает в остатке 2, …, на восьмую цифру –
8. Прав ли Знайка?

\problem{\textbf{6}} Прогулочный катер за неделю совершил 40 рейсов. На каждом рейсе было
ровно 10 пассажиров. Известно, что любые двое были вместе не более чем на
одном рейсе. Докажите, что общее количество пассажиров было не менее 60.